\chapter{Fundamento Teórico}


\section{Principio de Fermat}
El fenómeno se basa en el Principio de Fermat, el cual define la trayectoria de un rayo de luz entre dos puntos $A$ y $B$ como la curva $\gamma$ que minimiza la integral del camino óptico:

$$L = \int_{A}^{B} n[\vec{r}] ds$$

En este caso, el índice de refracción $n[\vec{r}]$ no es constante, sino que se define como un campo aleatorio con correlación espacial:

$$n[\vec{r}] = n_0 + \delta n[\vec{r}]$$

Aquí $\delta n[\vec{r}]$ representa las fluctuaciones de densidad. La acumulación de pequeñas desviaciones debidas a estas variaciones provoca el enfoque y desenfoque de los rayos, resultando en la formación del flujo ramificado.


\section{Campos Aleatorios y Desorden Correlacionado}
Para modelar el medio de densidad variable de manera rigurosa, las fluctuaciones espaciales se tratan matemáticamente mediante procesos estocásticos. El potencial del entorno se representa típicamente mediante campos aleatorios gaussianos. Para garantizar la aparición del flujo ramificado, la longitud de correlación de este desorden debe ser mayor que la longitud de onda de la luz incidente, y la amplitud de las variaciones debe ser pequeña. Estas condiciones aseguran que la dispersión ocurra de forma suave y predominantemente hacia adelante.

\section{Dinámica Hamiltoniana}
La formulación geométrica obtenida del Principio de Fermat puede traducirse al lenguaje de la mecánica analítica. La propagación de los rayos a través del medio se rige por las ecuaciones de Hamilton. Durante su recorrido por el potencial estocástico, las continuas refracciones inducen cambios en el momento transversal de cada rayo. Este formalismo permite calcular de manera precisa la evolución del sistema a medida que las perturbaciones se acumulan a lo largo del camino óptico.

\section{Formación de Cáusticas}
La acumulación de las desviaciones transversales altera la distribución espacial de los rayos, provocando que el espacio de fase se pliegue sobre sí mismo. Al proyectar la dinámica del espacio de fase en el espacio de coordenadas físicas, se generan singularidades donde la densidad de rayos es extremadamente alta. Estas regiones de máxima intensidad luminosa se conocen como cáusticas y corresponden a los filamentos principales que estructuran visualmente el flujo.

\section{Análisis en Perspectiva de Redes}
La compleja topología que generan las cáusticas al propagarse proporciona el marco adecuado para aplicar la teoría de redes complejas. Bajo este esquema, los puntos donde los filamentos de luz se dividen o se cruzan actúan como nodos, mientras que los canales de alta intensidad luminosa funcionan como enlaces. La construcción de grafos a partir de estos elementos permite evaluar el flujo ramificado de la luz en medios de densidad variable analizando el transporte de energía y las propiedades estadísticas del sistema mediante métricas topológicas.